\notechapter{N-20220412}{Linear Algebra Continued: Indices, Rank, Transpose, Hermitian Conjugate, and Trace}

\section{Unit vectors and Kronecker delta}

This note follows the previous linear algebra material and collects notation that will be useful later.

\begin{definition}
A unit vector is a vector of length \(1\).  I write such a vector as \(\mathbf v\) when the direction matters more than its magnitude.
\end{definition}

\begin{definition}[Kronecker delta]
The Kronecker delta is
\[
  \delta_{ij}=\begin{cases}
  1,&i=j,\\
  0,&i\neq j.
  \end{cases}
\]
Thus the matrix \((\delta_{ij})\) is the identity matrix:
\[
  (\delta_{ij})=
  \begin{pmatrix}
  1&0&0\\
  0&1&0\\
  0&0&1
  \end{pmatrix}.
\]
\end{definition}

\section{Einstein summation convention}

\begin{definition}
Einstein summation convention says that when an index is repeated exactly twice in a term, summation over that index is understood:
\[
  x_iy_i=\sum_i x_iy_i.
\]
A repeated index is called a dummy index, and its name does not matter:
\[
  x_iy_i=x_jy_j=x_ky_k.
\]
\end{definition}

The basic rules are:
\begin{enumerate}[(i)]
  \item an index appearing once is a free index;
  \item an index appearing twice is summed over;
  \item an index appearing three or more times in one term is not allowed.
\end{enumerate}

\section{Rank}

\begin{definition}
The column rank of a matrix is the maximal number of linearly independent columns.  The row rank is the maximal number of linearly independent rows.  These two numbers are equal, and their common value is the rank of the matrix.
\end{definition}

Gaussian elimination is the standard way to compute rank: reduce the matrix to row echelon form and count the nonzero rows.


\section{Transpose and Hermitian conjugate}

\begin{definition}
If \(A\) is an \(m\times n\) matrix, its transpose \(A^T\) is the \(n\times m\) matrix defined by
\[
  (A^T)_{ij}=A_{ji}.
\]
\end{definition}

\begin{definition}
The Hermitian conjugate of \(A\) is
\[
  A^*=(A^T)^{\overline{\phantom{x}}}=\overline{A}^{\,T}.
\]
It satisfies
\[
  (AB)^*=B^*A^*.
\]
Indeed,
\[
  (AB)^*=\overline{(AB)}^T=(\overline A\,\overline B)^T=\overline B^T\overline A^T=B^*A^*.
\]
\end{definition}

\begin{definition}
A matrix is symmetric if
\[
  A^T=A.
\]
A matrix is Hermitian if
\[
  A^*=A.
\]
A matrix is skew-Hermitian if
\[
  A^*=-A.
\]
In a skew-Hermitian matrix, the diagonal entries are purely imaginary.
\end{definition}


\section{Trace}

\begin{definition}
The trace of an \(n\times n\) matrix \(A\) is the sum of its diagonal entries:
\[
  \operatorname{tr}(A)=\sum_i A_{ii}.
\]
In Einstein notation this is simply
\[
  \operatorname{tr}(A)=A_{ii}.
\]
\end{definition}


\section{A more structural view}

The common theme is index notation.  The Kronecker delta represents the identity, repeated indices encode summation, and matrix operations such as transpose, adjoint, and trace become compact when written with indices.

\section{Transpose and dual maps}

If $A:V\to W$ is a linear map, the dual map is
\[
  A^*:W^*\to V^*,\qquad A^*(\varphi)=\varphi\circ A.
\]
In bases, this dual map is represented by the transpose matrix.  This explains why transposition reverses order:
\[
  (AB)^T=B^TA^T.
\]
Composition of maps reverses when pulled back to functionals.

The elementary transpose is therefore not just flipping a matrix; it is the coordinate form of passing from vectors to covectors.

\section{Hermitian adjoints in Hilbert spaces}

In an inner product space, the adjoint $A^*$ is defined by
\[
  \langle Ax,y\rangle=\langle x,A^*y\rangle.
\]
Over $\mathbb C$, its matrix is the conjugate transpose.  A Hermitian matrix satisfies $A=A^*$ and has real eigenvalues.  A unitary matrix satisfies $A^*A=I$ and preserves the inner product.

The elementary real symmetric matrix is thus the real finite-dimensional case of self-adjoint operators.  This becomes central in quantum mechanics, Fourier analysis, and spectral theory.

\section{Trace as invariant}

The trace is the sum of diagonal entries:
\[
  \operatorname{tr}A=\sum_i a_{ii}.
\]
It is invariant under similarity because
\[
  \operatorname{tr}(P^{-1}AP)=\operatorname{tr}(APP^{-1})=\operatorname{tr}A.
\]
The cyclic rule
\[
  \operatorname{tr}(AB)=\operatorname{tr}(BA)
\]
is the key identity.

In category-theoretic language, trace abstracts the idea of closing an input-output wire.  In finite-dimensional vector spaces, that abstract trace reduces to the familiar matrix trace.  This is why traces appear in representation theory, statistical mechanics, and index theory.

\section{Schatten norms}

Rank counts nonzero singular values, while trace and determinant combine eigenvalues.  A unified family of matrix norms uses singular values:
\[
  \|A\|_{S_p}=\left(\sum_i \sigma_i(A)^p\right)^{1/p}.
\]
For $p=2$ this is the Frobenius norm; for $p=1$ it is the nuclear norm.  These norms are finite-dimensional models of operator ideals.

Thus rank, trace, Hermitian adjoint, and transpose are all pieces of a larger operator-theoretic vocabulary.

\section{Indices as tensor bookkeeping}

The Kronecker delta \(\delta_i^j\) is the coordinate symbol for the identity map.  Einstein summation records contractions, such as
\[
  y_i=A_i^{\ j}x_j.
\]
The repeated index is summed because a linear map consumes one coordinate direction and produces another.

\section{Transpose as pullback on dual spaces}

For \(A:V\to W\), the dual map is
\[
  A^*:W^*\to V^*,\qquad A^*(\varphi)=\varphi\circ A.
\]
In bases, this is represented by the transpose.  Hermitian conjugate adds complex conjugation because complex inner products are conjugate-linear in one argument.

\section{Trace as contraction}

The trace
\[
  \operatorname{tr}(A)=A_i^{\ i}
\]
is the contraction of one upper and one lower index.  It is invariant under change of basis because it is not really a coordinate sum; it is the categorical trace of an endomorphism in finite-dimensional vector spaces.

\section{Rank and factorization}

Rank is the dimension of the image.  A rank-\(r\) map factors through an \(r\)-dimensional space:
\[
  V\twoheadrightarrow \operatorname{im}A\hookrightarrow W.
\]
This explains why rank measures the true number of independent output directions.

\section{Coordinate indices as intrinsic linear algebra}

Indices, transpose, Hermitian conjugate, rank, and trace are coordinate shadows of intrinsic operations: duality, contraction, image factorization, and inner-product-dependent adjoints.
